% Конфиденциальное дополнение. Утверждения становятся глобальными только
% после отдельного доказательства, что каждый минимизатор ранга 10 содержит
% sqrt(3) E8. До этого K ограничивается двумя известными типами равенства.

\begin{proposition}[Внутренняя структура случаев равенства ранга десять]
\label{prop:intrinsic-rank-ten}
Пусть $K$ принадлежит одному из двух типов равенства ранга десять. Тогда:
\begin{enumerate}
\item $m(K^*)=2/5$, причём $K^*$ имеет ровно шесть минимальных векторов.
      Их линейная оболочка есть двумерная плоскость $V$, а порождённая
      ими решётка изометрична $A_2/\sqrt5$.
\item В $K$ имеется ровно одна подрешётка, изометричная $\sqrt3E_8$.
      Она определяется самой решёткой $K$ по формуле
      \[
        J=K\cap V^\perp.
      \]
\item В $K$ имеется ровно три примитивные подрешётки ранга девять
      определителя $21870$. Все три содержат $J$ и изометричны $K_9$.
\end{enumerate}
После глобальной классификации равенства предложение применимо ко всем
положительно определённым чётным целочисленным решёткам ранга десять
с минимумом не менее $6$ и определителем $54675$.
\end{proposition}

\begin{proof}
Выберем копию $J\cong\sqrt3E_8$ в $K$. Оценка радиуса классов
склейки $8/3<6$ обеспечивает примитивность $J$. Равенство в оценке
определителя бинарной решётки даёт для проекции $P$ на ортогональное
дополнение к $J$ матрицу Грама
\[
 \frac53\begin{pmatrix}2&1\\1&2\end{pmatrix}.
\]
Поэтому перпендикулярная двойственная решётка
\[
 T=K^*\cap\operatorname{span}(J)^\perp
\]
имеет матрицу Грама
\[
 \frac15\begin{pmatrix}2&-1\\-1&2\end{pmatrix}.
\]
У неё ровно шесть векторов нормы $2/5$, порождающих $T$.
Проекция любого вектора $K^*$ на $\operatorname{span}(J)$ принадлежит
$J^*$. Каждый ненулевой вектор $J^*\cong E_8/\sqrt3$ имеет норму
не менее $2/3$. Следовательно, все векторы $K^*$ нормы не более
$2/5$ лежат в $T$. Это доказывает первый пункт и определяет $V$
независимо от первоначального выбора копии $J$.

Применим тот же аргумент к любой другой копии $J'$ решётки
$\sqrt3E_8$. Её перпендикулярная двойственная решётка содержит те же
шесть минимальных векторов, порождающих плоскость $V$. Поэтому
$\operatorname{span}(J')=V^\perp$. Автоматическая примитивность даёт
$J'=K\cap V^\perp=J$.

Каждая примитивная подрешётка $H\subset K$ ранга девять является ядром
примитивного вектора $y\in K^*$, определённого с точностью до знака:
\[
 H=\{x\in K:\langle x,y\rangle=0\}.
\]
Здесь примитивность $y$ означает, что отображение
$x\mapsto\langle x,y\rangle$ переводит $K$ на $\mathbb Z$.
Прообраз единицы имеет ортогональную высоту $1/\|y\|$, поэтому
формула коволюмов даёт
\[
 \det H=(\det K)\|y\|^2.
\]
Следовательно, $\det H=21870$ тогда и только тогда, когда
$\|y\|^2=2/5$. Все шесть минимальных двойственных векторов
примитивны, а различные пары $\{y,-y\}$ дают различные ядра.
Получаются ровно три ядра. Поскольку $y\in V$, все они содержат $J$.
Их одномерная проекция перпендикулярно $J$ имеет квадрат шага
$21870/6561=10/3$. Поэтому класс склейки обязан быть глубоким
классом $J^*/J$ с минимумом $8/3$. Транзитивность на всех $1920$
глубоких классах показывает, что каждое такое расширение ранга девять
изометрично $K_9$.
\end{proof}

В предъявленных базисах шесть примитивных столбцов скалярных произведений
двойственных векторов имеют вид
\[
 (0,\ldots,0,\pm1,0),\quad
 (0,\ldots,0,0,\pm1),\quad
 (0,\ldots,0,\pm1,\pm1),
\]
где в последней паре знаки совпадают. Базисы трёх ядер получаются
добавлением к восьми базисным векторам $J$ соответственно векторов
$e_{10}$, $e_9$ и $e_9-e_{10}$. Последний из этих образующих сам
может не быть кратчайшим; сдвиг на вектор $J$ даёт представитель нормы $6$.

\begin{corollary}[Равномерная потеря плотности для нетривиальных кластеров]
\label{cor:uniform-rank-ten-cluster-loss}
Пусть $K$ принадлежит одному из двух типов равенства ранга десять,
а $S\subset K^*$ --- конечное множество с $|S|>1$ и квадратом
диаметра $\beta<6$. При гарантированном радиусе антиподной конструкции
$\sqrt{6-\beta}/2$ в размерности $38$ отношение к одноточечной
конструкции над тем же сечением при радиусе $\sqrt6/2$ удовлетворяет
\[
 \frac{\delta(S,K)}{\delta(\{0\},K)}
 =|S|(1-\beta/6)^{19}
 \le 9\left(\frac89\right)^{19}
 =0.9601621523027026\ldots<1.
\]
Таким образом, каждый нетривиальный кластер теряет не менее
$3.9837847697297\ldots\%$ этой гарантированной центральной плотности.
Константа является доказанной верхней оценкой; её точность не утверждается.
После глобальной классификации равенства результат охватывает все
сечения ранга десять минимального определителя.
\end{corollary}

\begin{proof}
Для каждого из десяти интервалов доказанной таблицы кластерных оценок
обозначим через $b$ включённый левый конец, а через $B$ --- оценку
числа точек. Используются пары
\[
 (2/5,3),\ (2/3,9),\ (23/30,11),\ (29/30,27),\ (6/5,36),
\]
\[
 (4/3,82),\ (8/5,109),\ (49/30,385),\ (2,1540),\ (8/3,54675).
\]
На каждом интервале монотонность даёт
$|S|(1-\beta/6)^{19}\le B(1-b/6)^{19}$.
Точное сравнение этих десяти рациональных чисел показывает, что
наибольшее равно $9(8/9)^{19}$ и соответствует паре $(2/3,9)$.
При $\beta<2/5$ нетривиальные кластеры отсутствуют.
\end{proof}

\begin{corollary}[Точная область вывода для размерности тридцать восемь]
\label{cor:rank-ten-density-scope}
Пусть $M$ --- положительно определённая чётная унимодулярная решётка
ранга $48$ с минимумом не менее $6$, а $K\subset M$ --- примитивное
сечение ранга десять. При радиусе $\sqrt6/2$ одноточечная конструкция
на ортогональном дополнении $N=M\cap K^\perp$ имеет центральную плотность
\[
 \delta_0(K)=\frac{(3/2)^{19}}{\sqrt{\det K}}
 \le\frac{(3/2)^{19}}{\sqrt{54675}}.
\]
Равенство в этом сравнении при гарантированном радиусе имеет место
ровно при $\det K=54675$. После глобальной классификации равенства
каждое такое сечение имеет описанную выше внутреннюю структуру,
а каждый нетривиальный кластер над ним испытывает доказанную
равномерную потерю плотности.
\end{corollary}

\begin{proof}
Примитивность и унимодулярность дают $\det N=\det K$.
Минимум $N$ не меньше $6$, а ранг равен $38$, поэтому выбранный радиус
даёт указанную центральную плотность. Теорема об определителе даёт
неравенство. Для конечного кластера обычное соответствие склейки
в унимодулярной решётке даёт один сдвиг $N$ на каждую точку кластера.
При $\beta<6$ эти сдвиги различны по модулю $N$: две точки кластера
из одного класса по модулю $K$ различались бы на ненулевой вектор $K$
нормы не менее $6$. Минимум объемлющей решётки $6$ обеспечивает
квадрат расстояния между различными сдвигами не меньше $6-\beta$.
Отсюда следует формула при гарантированном радиусе.
\end{proof}

Эти утверждения не ограничивают фактическую плотность, получаемую
увеличением радиуса, если конкретное ортогональное дополнение имеет
минимум больше $6$ или конкретный кластер даёт лучшее разделение,
чем $6-\beta$. Они не сравнивают нетривиальные кластеры над сечениями
большего определителя с оптимальной одноточечной конструкцией
и не доказывают оптимальность среди всех упаковок размерности $38$.

Для вложений в объемлющие решётки структурное предложение показывает,
что каждое оптимальное сечение ранга десять содержит ровно три копии
$K_9$. Поэтому глобальную классификацию таких вложений можно строить
через расширения вложений $K_9$. Уже проверенный каталог из $19$
расширений относится только к фиксированной напечатанной копии
$K_9$ в напечатанной решётке $P_{48p}$; он не является ни каталогом
всех оптимальных сечений ранга десять в этой решётке, ни подсчётом
орбит её группы автоморфизмов.

\paragraph{Независимый точный аудит.}
Скрипт \texttt{audit\_two\_types\_and\_cluster\_gap.py}, использующий
только стандартную библиотеку, имеет собственный рациональный перебор
эллипсоида и собственную матричную арифметику. Он проверяет все $6561$
классов вычетов, транзитивность на $1920$ глубоких классах, орбиты
второго класса размеров $560+1$ и $559$ явных целочисленных
унимодулярных конгруэнций матриц Грама размера $10\times10$, связывающих
большую орбиту. Ранги двух матриц Грама по модулю $3$ равны $4$ и $3$,
поэтому эти типы не могут стать изометричными после забывания выбранной
копии $E_8$. Прямой перебор в двойственных решётках находит ровно шесть
векторов нормы $2/5$ у каждого типа. Каждая из трёх матриц ядер имеет
определитель $21870$, минимум $6$ и $258$ векторов нормы $6$.
Наконец, все десять сравнений для равномерной кластерной оценки
проверяются как точные рациональные неравенства.
